% concatenated from arxiv-v2.tex, math_commands.tex
\documentclass{article} 

\usepackage{natbib}
%\usepackage[top=2.5cm,bottom=3cm,left=4cm,right=4cm]{geometry}
\usepackage[top=2.9cm,bottom=3.4cm,left=4cm,right=4cm]{geometry}

\renewcommand{\topfraction}{0.95} 
\renewcommand{\textfraction}{0.05}

%%%%% NEW MATH DEFINITIONS %%%%%

\usepackage{amsmath,amsfonts,bm}

% Mark sections of captions for referring to divisions of figures
\newcommand{\figleft}{{\em (Left)}}
\newcommand{\figcenter}{{\em (Center)}}
\newcommand{\figright}{{\em (Right)}}
\newcommand{\figtop}{{\em (Top)}}
\newcommand{\figbottom}{{\em (Bottom)}}
\newcommand{\captiona}{{\em (a)}}
\newcommand{\captionb}{{\em (b)}}
\newcommand{\captionc}{{\em (c)}}
\newcommand{\captiond}{{\em (d)}}

% Highlight a newly defined term
\newcommand{\newterm}[1]{{\bf #1}}


% Figure reference, lower-case.
\def\figref#1{figure~\ref{#1}}
% Figure reference, capital. For start of sentence
\def\Figref#1{Figure~\ref{#1}}
\def\twofigref#1#2{figures \ref{#1} and \ref{#2}}
\def\quadfigref#1#2#3#4{figures \ref{#1}, \ref{#2}, \ref{#3} and \ref{#4}}
% Section reference, lower-case.
\def\secref#1{section~\ref{#1}}
% Section reference, capital.
\def\Secref#1{Section~\ref{#1}}
% Reference to two sections.
\def\twosecrefs#1#2{sections \ref{#1} and \ref{#2}}
% Reference to three sections.
\def\secrefs#1#2#3{sections \ref{#1}, \ref{#2} and \ref{#3}}
% Reference to an equation, lower-case.
\def\eqref#1{equation~\ref{#1}}
% Reference to an equation, upper case
\def\Eqref#1{Equation~\ref{#1}}
% A raw reference to an equation---avoid using if possible
\def\plaineqref#1{\ref{#1}}
% Reference to a chapter, lower-case.
\def\chapref#1{chapter~\ref{#1}}
% Reference to an equation, upper case.
\def\Chapref#1{Chapter~\ref{#1}}
% Reference to a range of chapters
\def\rangechapref#1#2{chapters\ref{#1}--\ref{#2}}
% Reference to an algorithm, lower-case.
\def\algref#1{algorithm~\ref{#1}}
% Reference to an algorithm, upper case.
\def\Algref#1{Algorithm~\ref{#1}}
\def\twoalgref#1#2{algorithms \ref{#1} and \ref{#2}}
\def\Twoalgref#1#2{Algorithms \ref{#1} and \ref{#2}}
% Reference to a part, lower case
\def\partref#1{part~\ref{#1}}
% Reference to a part, upper case
\def\Partref#1{Part~\ref{#1}}
\def\twopartref#1#2{parts \ref{#1} and \ref{#2}}

\def\ceil#1{\lceil #1 \rceil}
\def\floor#1{\lfloor #1 \rfloor}
\def\1{\bm{1}}
\newcommand{\train}{\mathcal{D}}
\newcommand{\valid}{\mathcal{D_{\mathrm{valid}}}}
\newcommand{\test}{\mathcal{D_{\mathrm{test}}}}

\def\eps{{\epsilon}}


% Random variables
\def\reta{{\textnormal{$\eta$}}}
\def\ra{{\textnormal{a}}}
\def\rb{{\textnormal{b}}}
\def\rc{{\textnormal{c}}}
\def\rd{{\textnormal{d}}}
\def\re{{\textnormal{e}}}
\def\rf{{\textnormal{f}}}
\def\rg{{\textnormal{g}}}
\def\rh{{\textnormal{h}}}
\def\ri{{\textnormal{i}}}
\def\rj{{\textnormal{j}}}
\def\rk{{\textnormal{k}}}
\def\rl{{\textnormal{l}}}
% rm is already a command, just don't name any random variables m
\def\rn{{\textnormal{n}}}
\def\ro{{\textnormal{o}}}
\def\rp{{\textnormal{p}}}
\def\rq{{\textnormal{q}}}
\def\rr{{\textnormal{r}}}
\def\rs{{\textnormal{s}}}
\def\rt{{\textnormal{t}}}
\def\ru{{\textnormal{u}}}
\def\rv{{\textnormal{v}}}
\def\rw{{\textnormal{w}}}
\def\rx{{\textnormal{x}}}
\def\ry{{\textnormal{y}}}
\def\rz{{\textnormal{z}}}

% Random vectors
\def\rvepsilon{{\mathbf{\epsilon}}}
\def\rvtheta{{\mathbf{\theta}}}
\def\rva{{\mathbf{a}}}
\def\rvb{{\mathbf{b}}}
\def\rvc{{\mathbf{c}}}
\def\rvd{{\mathbf{d}}}
\def\rve{{\mathbf{e}}}
\def\rvf{{\mathbf{f}}}
\def\rvg{{\mathbf{g}}}
\def\rvh{{\mathbf{h}}}
\def\rvu{{\mathbf{i}}}
\def\rvj{{\mathbf{j}}}
\def\rvk{{\mathbf{k}}}
\def\rvl{{\mathbf{l}}}
\def\rvm{{\mathbf{m}}}
\def\rvn{{\mathbf{n}}}
\def\rvo{{\mathbf{o}}}
\def\rvp{{\mathbf{p}}}
\def\rvq{{\mathbf{q}}}
\def\rvr{{\mathbf{r}}}
\def\rvs{{\mathbf{s}}}
\def\rvt{{\mathbf{t}}}
\def\rvu{{\mathbf{u}}}
\def\rvv{{\mathbf{v}}}
\def\rvw{{\mathbf{w}}}
\def\rvx{{\mathbf{x}}}
\def\rvy{{\mathbf{y}}}
\def\rvz{{\mathbf{z}}}

% Elements of random vectors
\def\erva{{\textnormal{a}}}
\def\ervb{{\textnormal{b}}}
\def\ervc{{\textnormal{c}}}
\def\ervd{{\textnormal{d}}}
\def\erve{{\textnormal{e}}}
\def\ervf{{\textnormal{f}}}
\def\ervg{{\textnormal{g}}}
\def\ervh{{\textnormal{h}}}
\def\ervi{{\textnormal{i}}}
\def\ervj{{\textnormal{j}}}
\def\ervk{{\textnormal{k}}}
\def\ervl{{\textnormal{l}}}
\def\ervm{{\textnormal{m}}}
\def\ervn{{\textnormal{n}}}
\def\ervo{{\textnormal{o}}}
\def\ervp{{\textnormal{p}}}
\def\ervq{{\textnormal{q}}}
\def\ervr{{\textnormal{r}}}
\def\ervs{{\textnormal{s}}}
\def\ervt{{\textnormal{t}}}
\def\ervu{{\textnormal{u}}}
\def\ervv{{\textnormal{v}}}
\def\ervw{{\textnormal{w}}}
\def\ervx{{\textnormal{x}}}
\def\ervy{{\textnormal{y}}}
\def\ervz{{\textnormal{z}}}

% Random matrices
\def\rmA{{\mathbf{A}}}
\def\rmB{{\mathbf{B}}}
\def\rmC{{\mathbf{C}}}
\def\rmD{{\mathbf{D}}}
\def\rmE{{\mathbf{E}}}
\def\rmF{{\mathbf{F}}}
\def\rmG{{\mathbf{G}}}
\def\rmH{{\mathbf{H}}}
\def\rmI{{\mathbf{I}}}
\def\rmJ{{\mathbf{J}}}
\def\rmK{{\mathbf{K}}}
\def\rmL{{\mathbf{L}}}
\def\rmM{{\mathbf{M}}}
\def\rmN{{\mathbf{N}}}
\def\rmO{{\mathbf{O}}}
\def\rmP{{\mathbf{P}}}
\def\rmQ{{\mathbf{Q}}}
\def\rmR{{\mathbf{R}}}
\def\rmS{{\mathbf{S}}}
\def\rmT{{\mathbf{T}}}
\def\rmU{{\mathbf{U}}}
\def\rmV{{\mathbf{V}}}
\def\rmW{{\mathbf{W}}}
\def\rmX{{\mathbf{X}}}
\def\rmY{{\mathbf{Y}}}
\def\rmZ{{\mathbf{Z}}}

% Elements of random matrices
\def\ermA{{\textnormal{A}}}
\def\ermB{{\textnormal{B}}}
\def\ermC{{\textnormal{C}}}
\def\ermD{{\textnormal{D}}}
\def\ermE{{\textnormal{E}}}
\def\ermF{{\textnormal{F}}}
\def\ermG{{\textnormal{G}}}
\def\ermH{{\textnormal{H}}}
\def\ermI{{\textnormal{I}}}
\def\ermJ{{\textnormal{J}}}
\def\ermK{{\textnormal{K}}}
\def\ermL{{\textnormal{L}}}
\def\ermM{{\textnormal{M}}}
\def\ermN{{\textnormal{N}}}
\def\ermO{{\textnormal{O}}}
\def\ermP{{\textnormal{P}}}
\def\ermQ{{\textnormal{Q}}}
\def\ermR{{\textnormal{R}}}
\def\ermS{{\textnormal{S}}}
\def\ermT{{\textnormal{T}}}
\def\ermU{{\textnormal{U}}}
\def\ermV{{\textnormal{V}}}
\def\ermW{{\textnormal{W}}}
\def\ermX{{\textnormal{X}}}
\def\ermY{{\textnormal{Y}}}
\def\ermZ{{\textnormal{Z}}}

% Vectors
\def\vzero{{\bm{0}}}
\def\vone{{\bm{1}}}
\def\vmu{{\bm{\mu}}}
\def\vtheta{{\bm{\theta}}}
\def\va{{\bm{a}}}
\def\vb{{\bm{b}}}
\def\vc{{\bm{c}}}
\def\vd{{\bm{d}}}
\def\ve{{\bm{e}}}
\def\vf{{\bm{f}}}
\def\vg{{\bm{g}}}
\def\vh{{\bm{h}}}
\def\vi{{\bm{i}}}
\def\vj{{\bm{j}}}
\def\vk{{\bm{k}}}
\def\vl{{\bm{l}}}
\def\vm{{\bm{m}}}
\def\vn{{\bm{n}}}
\def\vo{{\bm{o}}}
\def\vp{{\bm{p}}}
\def\vq{{\bm{q}}}
\def\vr{{\bm{r}}}
\def\vs{{\bm{s}}}
\def\vt{{\bm{t}}}
\def\vu{{\bm{u}}}
\def\vv{{\bm{v}}}
\def\vw{{\bm{w}}}
\def\vx{{\bm{x}}}
\def\vy{{\bm{y}}}
\def\vz{{\bm{z}}}

% Elements of vectors
\def\evalpha{{\alpha}}
\def\evbeta{{\beta}}
\def\evepsilon{{\epsilon}}
\def\evlambda{{\lambda}}
\def\evomega{{\omega}}
\def\evmu{{\mu}}
\def\evpsi{{\psi}}
\def\evsigma{{\sigma}}
\def\evtheta{{\theta}}
\def\eva{{a}}
\def\evb{{b}}
\def\evc{{c}}
\def\evd{{d}}
\def\eve{{e}}
\def\evf{{f}}
\def\evg{{g}}
\def\evh{{h}}
\def\evi{{i}}
\def\evj{{j}}
\def\evk{{k}}
\def\evl{{l}}
\def\evm{{m}}
\def\evn{{n}}
\def\evo{{o}}
\def\evp{{p}}
\def\evq{{q}}
\def\evr{{r}}
\def\evs{{s}}
\def\evt{{t}}
\def\evu{{u}}
\def\evv{{v}}
\def\evw{{w}}
\def\evx{{x}}
\def\evy{{y}}
\def\evz{{z}}

% Matrix
\def\mA{{\bm{A}}}
\def\mB{{\bm{B}}}
\def\mC{{\bm{C}}}
\def\mD{{\bm{D}}}
\def\mE{{\bm{E}}}
\def\mF{{\bm{F}}}
\def\mG{{\bm{G}}}
\def\mH{{\bm{H}}}
\def\mI{{\bm{I}}}
\def\mJ{{\bm{J}}}
\def\mK{{\bm{K}}}
\def\mL{{\bm{L}}}
\def\mM{{\bm{M}}}
\def\mN{{\bm{N}}}
\def\mO{{\bm{O}}}
\def\mP{{\bm{P}}}
\def\mQ{{\bm{Q}}}
\def\mR{{\bm{R}}}
\def\mS{{\bm{S}}}
\def\mT{{\bm{T}}}
\def\mU{{\bm{U}}}
\def\mV{{\bm{V}}}
\def\mW{{\bm{W}}}
\def\mX{{\bm{X}}}
\def\mY{{\bm{Y}}}
\def\mZ{{\bm{Z}}}
\def\mBeta{{\bm{\beta}}}
\def\mPhi{{\bm{\Phi}}}
\def\mLambda{{\bm{\Lambda}}}
\def\mSigma{{\bm{\Sigma}}}

% Tensor
\DeclareMathAlphabet{\mathsfit}{\encodingdefault}{\sfdefault}{m}{sl}
\SetMathAlphabet{\mathsfit}{bold}{\encodingdefault}{\sfdefault}{bx}{n}
\newcommand{\tens}[1]{\bm{\mathsfit{#1}}}
\def\tA{{\tens{A}}}
\def\tB{{\tens{B}}}
\def\tC{{\tens{C}}}
\def\tD{{\tens{D}}}
\def\tE{{\tens{E}}}
\def\tF{{\tens{F}}}
\def\tG{{\tens{G}}}
\def\tH{{\tens{H}}}
\def\tI{{\tens{I}}}
\def\tJ{{\tens{J}}}
\def\tK{{\tens{K}}}
\def\tL{{\tens{L}}}
\def\tM{{\tens{M}}}
\def\tN{{\tens{N}}}
\def\tO{{\tens{O}}}
\def\tP{{\tens{P}}}
\def\tQ{{\tens{Q}}}
\def\tR{{\tens{R}}}
\def\tS{{\tens{S}}}
\def\tT{{\tens{T}}}
\def\tU{{\tens{U}}}
\def\tV{{\tens{V}}}
\def\tW{{\tens{W}}}
\def\tX{{\tens{X}}}
\def\tY{{\tens{Y}}}
\def\tZ{{\tens{Z}}}


% Graph
\def\gA{{\mathcal{A}}}
\def\gB{{\mathcal{B}}}
\def\gC{{\mathcal{C}}}
\def\gD{{\mathcal{D}}}
\def\gE{{\mathcal{E}}}
\def\gF{{\mathcal{F}}}
\def\gG{{\mathcal{G}}}
\def\gH{{\mathcal{H}}}
\def\gI{{\mathcal{I}}}
\def\gJ{{\mathcal{J}}}
\def\gK{{\mathcal{K}}}
\def\gL{{\mathcal{L}}}
\def\gM{{\mathcal{M}}}
\def\gN{{\mathcal{N}}}
\def\gO{{\mathcal{O}}}
\def\gP{{\mathcal{P}}}
\def\gQ{{\mathcal{Q}}}
\def\gR{{\mathcal{R}}}
\def\gS{{\mathcal{S}}}
\def\gT{{\mathcal{T}}}
\def\gU{{\mathcal{U}}}
\def\gV{{\mathcal{V}}}
\def\gW{{\mathcal{W}}}
\def\gX{{\mathcal{X}}}
\def\gY{{\mathcal{Y}}}
\def\gZ{{\mathcal{Z}}}

% Sets
\def\sA{{\mathbb{A}}}
\def\sB{{\mathbb{B}}}
\def\sC{{\mathbb{C}}}
\def\sD{{\mathbb{D}}}
% Don't use a set called E, because this would be the same as our symbol
% for expectation.
\def\sF{{\mathbb{F}}}
\def\sG{{\mathbb{G}}}
\def\sH{{\mathbb{H}}}
\def\sI{{\mathbb{I}}}
\def\sJ{{\mathbb{J}}}
\def\sK{{\mathbb{K}}}
\def\sL{{\mathbb{L}}}
\def\sM{{\mathbb{M}}}
\def\sN{{\mathbb{N}}}
\def\sO{{\mathbb{O}}}
\def\sP{{\mathbb{P}}}
\def\sQ{{\mathbb{Q}}}
\def\sR{{\mathbb{R}}}
\def\sS{{\mathbb{S}}}
\def\sT{{\mathbb{T}}}
\def\sU{{\mathbb{U}}}
\def\sV{{\mathbb{V}}}
\def\sW{{\mathbb{W}}}
\def\sX{{\mathbb{X}}}
\def\sY{{\mathbb{Y}}}
\def\sZ{{\mathbb{Z}}}

% Entries of a matrix
\def\emLambda{{\Lambda}}
\def\emA{{A}}
\def\emB{{B}}
\def\emC{{C}}
\def\emD{{D}}
\def\emE{{E}}
\def\emF{{F}}
\def\emG{{G}}
\def\emH{{H}}
\def\emI{{I}}
\def\emJ{{J}}
\def\emK{{K}}
\def\emL{{L}}
\def\emM{{M}}
\def\emN{{N}}
\def\emO{{O}}
\def\emP{{P}}
\def\emQ{{Q}}
\def\emR{{R}}
\def\emS{{S}}
\def\emT{{T}}
\def\emU{{U}}
\def\emV{{V}}
\def\emW{{W}}
\def\emX{{X}}
\def\emY{{Y}}
\def\emZ{{Z}}
\def\emSigma{{\Sigma}}

% entries of a tensor
% Same font as tensor, without \bm wrapper
\newcommand{\etens}[1]{\mathsfit{#1}}
\def\etLambda{{\etens{\Lambda}}}
\def\etA{{\etens{A}}}
\def\etB{{\etens{B}}}
\def\etC{{\etens{C}}}
\def\etD{{\etens{D}}}
\def\etE{{\etens{E}}}
\def\etF{{\etens{F}}}
\def\etG{{\etens{G}}}
\def\etH{{\etens{H}}}
\def\etI{{\etens{I}}}
\def\etJ{{\etens{J}}}
\def\etK{{\etens{K}}}
\def\etL{{\etens{L}}}
\def\etM{{\etens{M}}}
\def\etN{{\etens{N}}}
\def\etO{{\etens{O}}}
\def\etP{{\etens{P}}}
\def\etQ{{\etens{Q}}}
\def\etR{{\etens{R}}}
\def\etS{{\etens{S}}}
\def\etT{{\etens{T}}}
\def\etU{{\etens{U}}}
\def\etV{{\etens{V}}}
\def\etW{{\etens{W}}}
\def\etX{{\etens{X}}}
\def\etY{{\etens{Y}}}
\def\etZ{{\etens{Z}}}

% The true underlying data generating distribution
\newcommand{\pdata}{p_{\rm{data}}}
% The empirical distribution defined by the training set
\newcommand{\ptrain}{\hat{p}_{\rm{data}}}
\newcommand{\Ptrain}{\hat{P}_{\rm{data}}}
% The model distribution
\newcommand{\pmodel}{p_{\rm{model}}}
\newcommand{\Pmodel}{P_{\rm{model}}}
\newcommand{\ptildemodel}{\tilde{p}_{\rm{model}}}
% Stochastic autoencoder distributions
\newcommand{\pencode}{p_{\rm{encoder}}}
\newcommand{\pdecode}{p_{\rm{decoder}}}
\newcommand{\precons}{p_{\rm{reconstruct}}}

\newcommand{\laplace}{\mathrm{Laplace}} % Laplace distribution

\newcommand{\E}{\mathbb{E}}
\newcommand{\Ls}{\mathcal{L}}
\newcommand{\R}{\mathbb{R}}
\newcommand{\emp}{\tilde{p}}
\newcommand{\lr}{\alpha}
\newcommand{\reg}{\lambda}
\newcommand{\rect}{\mathrm{rectifier}}
\newcommand{\softmax}{\mathrm{softmax}}
\newcommand{\sigmoid}{\sigma}
\newcommand{\softplus}{\zeta}
\newcommand{\KL}{D_{\mathrm{KL}}}
\newcommand{\Var}{\mathrm{Var}}
\newcommand{\standarderror}{\mathrm{SE}}
\newcommand{\Cov}{\mathrm{Cov}}
% Wolfram Mathworld says $L^2$ is for function spaces and $\ell^2$ is for vectors
% But then they seem to use $L^2$ for vectors throughout the site, and so does
% wikipedia.
\newcommand{\normlzero}{L^0}
\newcommand{\normlone}{L^1}
\newcommand{\normltwo}{L^2}
\newcommand{\normlp}{L^p}
\newcommand{\normmax}{L^\infty}

\newcommand{\parents}{Pa} % See usage in notation.tex. Chosen to match Daphne's book.

\DeclareMathOperator*{\argmax}{arg\,max}
\DeclareMathOperator*{\argmin}{arg\,min}

\DeclareMathOperator{\sign}{sign}
\DeclareMathOperator{\Tr}{Tr}
\let\ab\allowbreak


\usepackage{url}
\usepackage{amsmath}
\usepackage{amsfonts}
\usepackage{amssymb}
\usepackage{amsthm}
\usepackage{hyperref}
\usepackage{algorithmic}
\usepackage{algorithm}
%\usepackage[noend]{algpseudocode}
\usepackage{graphicx}
\usepackage{subcaption}
\usepackage{enumitem}
\usepackage{float}
\usepackage{xcolor} % Required for the color
\renewcommand{\algorithmiccomment}[1]{\hfill \textcolor{gray}{// #1}}

\theoremstyle{plain}
\newtheorem{theorem}{Theorem}[section]
\newtheorem{proposition}[theorem]{Proposition}
\newtheorem{lemma}[theorem]{Lemma}
\newtheorem{corollary}[theorem]{Corollary}
\theoremstyle{definition}
\newtheorem{definition}[theorem]{Definition}
\newtheorem{assumption}[theorem]{Assumption}
\theoremstyle{remark}
\newtheorem{remark}[theorem]{Remark}




%\newcommand{\R}{\mathbb{R}}
%\newcommand{\E}{\mathbb{E}}
\newcommand{\N}{\mathbb{N}}
\newcommand{\cN}{\mathcal{N}}
\renewcommand{\d}{\mathrm{d}}
\newcommand{\grad}[1]{\nabla_{\! #1}}
\newcommand{\ddgamma}{\tfrac{\d}{\d\gamma}}
\newcommand{\ddsigma}{\tfrac{\d}{\d\sigma}}
\newcommand{\ddxi}{\tfrac{\partial}{\partial x_i}}
\newcommand{\g}{g_\gamma}
\newcommand{\sumi}{\sum_{i\in[n]}}
\newcommand{\sumj}{\sum_{j\in[n]}}




\title{MAD: Manifold Attracted Diffusion}



\author{ Dennis Elbrächter\thanks{Department of Mathematics, University of Vienna,
Austria (dennis.elbraechter@univie.ac.at).} \and Giovanni S.~Alberti\footnotemark[2]\thanks{MaLGa Center, Department of Mathematics, University of Genoa, Italy (giovanni.alberti@unige.it, matteo.santacesaria@unige.it).} \and Matteo Santacesaria\footnotemark[2]}



\newcommand{\fix}{\marginpar{FIX}}
\newcommand{\new}{\marginpar{NEW}}
%\printAffiliationsAndNotice{\icmlEqualContribution}

\begin{document}
\maketitle


\begin{abstract}
Score-based diffusion models are a highly effective method for generating samples from a distribution of images. We consider scenarios where the training data comes from a noisy version of the target distribution, and present an efficiently implementable modification of the inference procedure to generate noiseless samples.
Our approach is motivated by the manifold hypothesis, according to which meaningful data is concentrated around some low-dimensional manifold of a high-dimensional ambient space. The central idea is that noise manifests as low magnitude variation in off-manifold directions in contrast to the relevant variation of the desired distribution which is mostly confined to on-manifold directions. We introduce the notion of an extended score and show that, in a simplified setting, it can be used to reduce small variations to zero, while leaving large variations mostly unchanged. We describe how its approximation can be computed efficiently from an approximation to the standard score and demonstrate its efficacy on toy problems, synthetic data, and real data. 
\end{abstract}

\section{Introduction}



%Generative models have become a cornerstone of modern machine learning, enabling the synthesis of high-fidelity data across domains such as images, audio, and text. 
Score-based diffusion models are a state-of-the-art approach for image synthesis, often outperforming alternatives like generative adversarial networks (GANs) and variational autoencoders (VAEs) in sample quality and diversity \citep{ho2020denoising,song2020score}. These models learn a neural network by adding noise to data samples during training, according to some forward process. The network can then be used to reverse this process during inference. Essentially, they denoise an image with independent Gaussian pixel values into samples from the data distribution.

However, real-world datasets are often corrupted by noise arising from measurement errors, compression artifacts, or data collection processes \citep{GUPTA2019466,Brummer_2019_CVPR_Workshops}. In this case, standard diffusion models trained directly on noisy data will learn to reproduce this corruption in their generated samples. We aim to address this by developing a method that can generate samples that  approximately come from the clean distribution, despite being trained on noisy data. 


%A key assumption underpinning the success of generative models, particularly in high-dimensional spaces like pixel data, is the \textit{manifold hypothesis}: 
A well-established paradigm for understanding data with a high-dimensional representation is the \textit{manifold hypothesis}: meaningful data distributions are concentrated near a low-dimensional manifold embedded in the ambient space \citep{bengio2013representation,fefferman2016testing}. For instance, natural images, despite being represented as a $d$-dimensional array, where $d$ can be in the millions, exhibit intrinsic dimensions that are far lower \citep{pope2021intrinsic}. 
%This hypothesis motivates techniques that exploit geometric structure to improve efficiency and robustness.
This viewpoint will be the basis of our approach as we interpret noise as low-magnitude variations in directions orthogonal to the manifold, whereas meaningful variations of the underlying data correspond to movements along the manifold itself. The central idea is to exploit this geometric structure during the inference process to suppress the former while preserving the latter.

The manifold hypothesis has previously been considered in the context of diffusion models to study their convergence behavior, e.g.\@ in \cite{debortoli2023convergence}, \cite{pmlr-v202-chen23o}, \cite{tang2024adaptivity}, and \cite{potaptchik2024linear}. It has also been argued in \cite{stanczuk2024diffusion} that diffusion models can be leveraged to estimate the intrinsic dimension of data manifolds.
Some existing works tackle the problem of noisy training \citep{daras2024ambient,lu2025stochastic}, but require adapting the training process. More precisely, the approach introduced in \cite{daras2024ambient} modifies the training loss to take into account the (unknown) degradation process, while the method in \cite{lu2025stochastic} is based on a specific training process where a small percentage of the data is clean. Outside the generative setting, there are also traditional manifold denoising techniques predating diffusion models \citep{hein2006manifold,gong2010manifold,wang2010manifold,fefferman2018fitting,JCAM-2023}. However, none of these approaches are efficient in generating new samples from the clean distribution.

This work introduces \textit{Manifold Attracted Diffusion (MAD)}, an efficiently implementable modification to the inference procedure of score-based models. We define the concept of an ``extended score'', which coincides with the standard score when that exists, but is also well-defined for Dirac delta distributions. In fact, it treats them essentially as Gaussians with a certain non-zero variance. This property can be leveraged in the inference procedure to reduce small off-manifold variations to almost zero while leaving larger on-manifold variations mostly unchanged, resulting in a soft thresholding effect. Thereby our method implicitly ``attracts'' the generated samples towards a low-dimensional structure, effectively filtering out noise.
Importantly, a suitable approximation to the extended score can be computed easily from an approximation to the standard score, which enables the use of established training methods as well as pretrained networks under some conditions that are, e.g., satisfied by the framework of \cite{edm}.
Our key contributions include:
\begin{itemize}[leftmargin=4mm]
    \item The formal definition and analysis of the extended score;
    \item An inference algorithm that reduces noise in generated samples without needing a special training procedure, making it compatible with established frameworks and pretrained models;
    \item Empirical validation on toy problems and real-world image data, such as FFHQ, AFHQ, ImageNet, and EMPIAR-11618 (cryo electron microscopy data).
\end{itemize}

We also distinguish our work from other diffusion-based approaches. Methods for posterior sampling, such as DPS \citep{chungdiffusion,chung2022improving}, are designed to solve inverse problems \citep{DBLP:journals/corr/abs-2410-00083}, such as denoising a single given image. Our goal is different: we learn from a dataset of noisy images to produce new, clean samples from the underlying distribution.

%Furthermore, unlike other inference-time quality improvements such as truncated sampling \citep{daras2025how}, MAD does not simply alter the sampling procedure. Instead, it introduces a novel, theoretically-grounded modification to the score function itself, using the extended score to enforce manifold adherence.


This paper is structured as follows. In Section~\ref{sec:background}, we review the necessary background on score-based diffusion models. In Section~\ref{sec:concept}, we formally introduce the extended score and analyze its properties. Section~\ref{sec:implementation} details our proposed inference algorithm and, in Section~\ref{sec:numerics}, we present numerical experiments that validate our approach. The related code is available at \url{https://github.com/delbraechter/MAD}.


% \subsection{Notation}
% [Not sure if necessary to have a subsection for this. Things to (potentially) define here: stochastic processes etc, normal distribution, convolution,... ]

\section{Background}\label{sec:background}
We will be working with the probability flow ODE formulation of diffusion models, largely following the framework of \cite{edm}. For more background and the connection to other diffusion model formulations we refer to \cite{edm} and the references therein. 

Given a data distribution $p_0$ on $\R^d$ we consider the stochastic process $(X_\sigma)_{\sigma\in[0,\sigma_{\rm max}]}$ with $X_0\sim p_0$ and
\begin{align*}
    X_\sigma = X_0 + \cN(0,\sigma^2\mathbb{I}),\quad \sigma>0.
\end{align*}
The corresponding densities are given by
\begin{align}\label{eq:psigma}
    p_\sigma = p_0 * g_{\sigma^2},
\end{align}
where $g_{\sigma^2}(x)=(2\pi \sigma^2)^{-d/2}\exp(-\frac{\|x\|^2}{2\sigma^2})$ is the density of $\cN(0,\sigma^2\mathbb{I})$. The central  idea behind score-based diffusion models is to generate a sample $x_0\sim X_0$ from a sample $x_{\sigma_{\rm max}} \sim X_{\sigma_{\rm max}}$, where $x_{\sigma_{\rm max}}$ is, in practice, approximated by simply sampling from $\cN(0,\sigma_{\rm max}^2\mathbb{I})$, as for a sufficiently large $\sigma_{\rm max}$ this should only introduce a negligible error. One way to achieve this is by evolving the ODE
\begin{align}\label{eq_pfODE}
    \d x_t = -\dot{\sigma}(t)\sigma(t) Sp_{\sigma(t)}(x_t)\d t,
\end{align}
where $\sigma\colon[0,T]\to[0,\sigma_{\rm max}]$ is some noise schedule and the score operator is given by
\begin{align*}
    S\colon P(\R^d)\to C(\R^d,\R^d),\qquad 
    p \mapsto \grad{x}\log p ,
\end{align*}
where $P(\R^d):=\{p\in C^1(\R^d,\R)\colon \int_{\R^d} p(x)\,dx=1,\; p(x)>0 \ \forall x\in\R^d\}$ is the set of densities that are positive everywhere in $\R^d$. 
Note that, for $\sigma>0$, $p_0 * g_{\sigma^2}$ is a density even if $p_0$ is not. The fact that the score may not be well-defined for $p_0$, e.g.\@ because the data distribution is supported on some lower-dimensional subset in $\R^d$, is avoided in practice by generating a sample $x_\delta$ with $\delta$ close to $0$.
It can be shown \citep{song2021,edm} that evolving a sample $x_{t_1}\sim X_{t_1}$ from $t_1$ to $t_2$ according to (\ref{eq_pfODE}) yields a sample $x_{t_2}\sim X_{t_2}$.

Of course, this is only useful provided that we have access to the score, which depends on the data distribution, and from which we usually  have only a finite number of samples. Remarkably, it turns out that a useful approximation of the score can be learned by training a neural network on these samples.
A common practice, motivated by numerical stability, is to learn the so-called denoiser function $D$, which is simply a shifted and scaled version of the score, i.e.\@ $D(x,\sigma) = \sigma^2 Sp_\sigma(x) + x$. This denoiser function is then approximated by a neural network $D_\theta$ using a loss based on
\begin{align*}
    \E_{y\sim X_0}\E_{\eta \sim \cN(0,\sigma^2\mathbb{I})}\|D_\theta(y + \eta)-y\|^2_2
\end{align*}
which is minimized by the denoiser function. However, as the expectation over the data distribution must be replaced by the empirical expectation based on the available data samples, this loss has, in general, many global minima. Nonetheless, employing neural networks with an adapted U-Net architecture \citep{unet1,ho2020denoising,unet2} appears to introduce sufficient bias towards a good approximation of the score. 

This approximation capability is primarily founded on the empirical observation that using it to generate images by evolving \eqref{eq_pfODE} (or related differential equations) produces realistic samples. An analytical description of how $D_\theta$ approximates $D$ is still an open problem. As the purpose of the work is to introduce a novel inference procedure, we will simply assume that we can obtain a suitable approximation to the score by some established training method.



\section{Concept: the Extended Score}\label{sec:concept}
Let $M(\R^d)$ denote the set of probability measures on $\R^d$.
\begin{definition}[Extended score]\label{def:exScore}
For $d \geq 1$, $p\in M(\R^d)$, $\gamma\in(0,\infty)$, let
\begin{align*}
     &H_\gamma\colon M(\R^d)\to C(\R^d,\R^d),\\    
     & p\mapsto (1+\gamma)S(p*g_\gamma) +\gamma\ddgamma S(p*g_\gamma)
     \end{align*}
     and
     \[
    H_0\colon \widetilde{M}(\R^d)\to \{f\colon \R^d\to\R^d\},\qquad 
    p\mapsto\lim_{\gamma\to 0} H_\gamma p,
    \]
where $$\widetilde{M}(\R^d):=\{p\in M(\R^d)\colon \lim_{\gamma\to 0} (H_\gamma p)(x)\in\R^d\ \forall x\in\R^d\}.$$
\end{definition}
We first note that $H_0$ coincides with the score for probability distributions with density in $P(\R^d)$, which we can view as a subset of $M(\R^d)$ by identifying a density function $p\in P(\R^d)$ with the measure given by $p(A)=\int_A p(x)\d x$ for $A\subseteq\R^d$ (see Section~\ref{sec:proofs} for a proof).
\begin{lemma}\label{lem:extension}
    Let $p\in P(\R^d)$. Then
    \(
        H_0p=Sp.
    \)
\end{lemma}
It is, however, also well-defined for, e.g., Dirac delta measures. Specifically, let $\delta$ denote the Dirac delta at $0$, then
\begin{align*}
    H_\gamma \delta(x)&=(1+\gamma)S(g_\gamma)(x)+\gamma\ddgamma S(g_\gamma)(x)\\&= -\tfrac{(1+\gamma)x}{\gamma} +\gamma\ddgamma(-\tfrac{x}{\gamma})=-x.
\end{align*}
In particular it holds that $H_0\delta(x)=-x$, i.e.\@ we obtain\footnote{Note that in this simple case, $H_\gamma\delta$ is already the same as $H_0\delta$, which is not the case in more complicated scenarios, e.g.\@ for mixtures of Dirac deltas.} a function which, similar to the score for Gaussians, yields at each point a vector pointing towards the mode of the probability distribution. This generalizes to mixtures of Dirac delta distributions, where $H_0p$ will point towards the location of the nearest Dirac delta in the mixture (see Section~\ref{sec:proofs} for a proof).\\
\begin{samepage}
\begin{lemma}\label{lem:Dirac_mix}
    Let $n\in\N$, $\mu_1,\dots,\mu_n\in\R^d$, $c_1,\dots,c_n\in\R_+$ such that $\sum_i c_i=1$ and 
   \(
        p=\sum_{i\in[n]} c_i\delta_{\mu_i}.
    \)
    Then
        \begin{align*}
        H_0p(x) = -\sum_{i\in[n]} z_i(x)(x-\mu_i),
    \end{align*}
    where $W_i=\{x\in\R^d\colon \|x-\mu_i\|\leq\|x-\mu_j\|\,\forall\,j\in[n]\}$ is a Voronoi region and
    \begin{align*}
        z_i(x)
        &=\begin{cases}
            0 & x\notin W_i,\\
            1 & x\in \operatorname{int} W_i,\\
            c_i\left(\sum_{j:x\in W_j} c_j\right)^{-1} & x\in \partial W_i.
        \end{cases}
    \end{align*} 
    \end{lemma}
\end{samepage}
    
The first two cases  in the expression for $z_i(x)$ cover almost every (w.r.t\@ the Lebesgue measure) $x\in\R^d$ and the third case is only needed if $x$ is equally distant to multiple $\mu_i$.
By combining the expressions of the extended score for a Dirac delta and for a non-degenerate Gaussian distribution, it is possible to derive an explicit expression of the extended score for any (possibly degenerate) Gaussian distribution. This is due to how the extended score behaves with respect to products of measures, see Appendix~\ref{sec:product} for the details.



We have the rather peculiar property that $H_0\delta = Sg_1$, i.e.\@ the extended score of the Dirac delta matches the score of a variance $1$ Gaussian\footnote{Changing the factor $(1+\gamma)$ to $(1+\tfrac{\gamma}{v})$ in the definition of $H_\gamma$ would lead to the equivalency $H_0\delta = Sg_v$ instead of $H_0\delta = Sg_1$. We will exploit the fact that this equivalency is adjustable in the next section, albeit only implicitly as we also need to account for the fact that we cannot take the $\gamma\to 0$ limit in practice.}. At this point, it should be noted that it is impossible to find a nice extension of the score operator that includes Dirac delta distributions, as
\begin{align*}
    Sg_\gamma(x)=-\tfrac{x}{\gamma},
\end{align*}
which diverges for $x\neq 0$ as $\gamma\to 0$. As such any extension is necessarily discontinuous with respect to any topology in which $\lim_{\gamma\to 0}g_\gamma=\delta$.

However, the ability of the extended score to treat distributions with positive variance like they were Dirac delta distributions will, in fact, be the cornerstone of our proposed inference technique. 
Under the manifold hypothesis, the clean  data distribution  has significant variation only along a small number of directions. Given noise in the ambient pixel space, which has a much higher dimension than the image manifold, the variance in off-manifold directions due to noise should be much smaller (e.g.\@ for isotropic Gaussian noise $\eta$ the variance in each direction is of order $\sim\! d^{-\frac{1}{2}} \E[\|\eta\|]$). Thus, noise can be suppressed by using the extended score. This principle is formally justified in Appendix~\ref{sec:product}, where we derive the extended score for a distribution supported on a low-dimensional subspace and show that it strongly attracts samples in the off-manifold directions while preserving the standard score in the on-manifold directions.

It is worth mentioning that the idea of decomposing the score function along the tangential and normal directions to the manifold was already exploited in 
 \cite{pmlr-v202-chen23o} (in the case of linear subspaces) and in \cite{liu2026improving} (in the case of general manifolds). Several interesting theoretical aspects are discussed in these papers, which, however, do not deal with the denoising task discussed in this work.



\section{Implementation}\label{sec:implementation}

%\gray{We would like to have an inference procedure which makes use of this property to generate samples which contain less of the noise than what would be obtained by running inference with the usual score.} 
In view of the property of the extended score just discussed, we would like to design an inference procedure that is able to generate samples with less noise, if compared to the samples obtained via the usual score. 
To this end we first note that, due to (\ref{eq:psigma}), we have
\begin{align}\begin{split}\label{eq:VE_trick}
S(p_{\sigma(t)}*g_\gamma)&=S(p_0*g_{\sigma(t)^2}*g_\gamma)\\
&=S(p_0*g_{\sigma(t)^2+\gamma})\\&=S(p_{\sqrt{\sigma(t)^2+\gamma}}).
\end{split}\end{align}
So, given a network trained to approximate the score for any $\sigma\in[0,\sigma_{\rm max}]$, as is the case in the framework of \cite{edm}, we also have an approximation to the extended score with a given small $\gamma>0$, where the $\gamma$-derivative can be obtained, e.g.\@, by a finite difference\footnote{Note that one could also compute the exact derivative of the score network using automatic differentiation (AD). However, since PyTorch is not optimized for forward mode AD this is significantly more expensive than a second evaluation of the network and did not yield any clear improvements in the numerical experiments.} approximation. As the desirable properties of the extended score hold in the $\gamma\to 0$ limit, which we cannot compute directly, we instead need to choose a suitable dependence $\gamma(t)$ with $\lim_{t\to 0}\gamma(t)=0$.

 A basic way to conduct inference with the standard score is evolving (\ref{eq_pfODE})
 simply via Euler method, i.e.\@ initializing with $x_0\sim\cN(0,\sigma(t_0)^2\mathbb{I})$ and iterating
\begin{align}\label{eq:basic_Euler}
    x_{i+1} = x_{i}- (t_{i+1}-t_i)\dot{\sigma}(t_i)\sigma(t_i) S_\theta(\sigma(t_i),x_i),
\end{align}
where $S_\theta(\sigma,x)\approx Sp_\sigma(x)$ is the learned approximation of the score. We will instead, based on Definition~\ref{def:exScore}, consider the iteration 
\small
\begin{equation}\label{eq:modified_inference}
    x_{i+1} = x_{i}- m(t_i)(t_{i+1}-t_i)\dot{\sigma}_\gamma(t_i)\sigma_\gamma(t_i) \alpha_i, 
\end{equation}
\normalsize
where
\[
\alpha_i=\big((1+\gamma(t_i))S_\theta(\sigma_\gamma(t_i),x_i) +\gamma(t_i)\ddgamma S_\theta(\sigma_\gamma(t_i),x_i)\big),\]
$\sigma_\gamma(t)=\sqrt{a\sigma(t)^2+b\gamma(t)}$, and $m(t_i)$ is a correction factor determined as explained below.

Here, $a,b>0$ are manually chosen parameters. Note that this reduces to (\ref{eq:basic_Euler}) for $a=1$, $b=0$, and $m(t)=(1+\gamma(t))^{-1}$, so essentially the choices of $a$, $b$, and $m$ determine to what extent we would like the inference to push points onto a manifold at a given time step. More precisely, increasing $a$ or reducing $b$ corresponds to lessening the effect of the extended score and thereby limiting the denoising effect. Conversely, reducing $a$ or increasing $b$ yield a stronger denoising effect. As such, $a$ and $b$ can be seen as regularization parameters, and their choice is problem dependent. While in the idealized case of the $\gamma\to 0$ limit, analyzed in Section~\ref{sec:concept}, Diracs and variance $1$ Gaussians are treated in the same way, the choice of $a$ and $b$ adjusts this equivalence in practice, causing variations below this (soft) threshold to be compressed significantly. Note that we currently do not determine the explicit dependence of this threshold on the parameters. However, since we cannot expect to know the optimal threshold for a data distribution in practice, a parameter optimization seems unavoidable anyway. We further enforce $\sigma_\gamma(t)=t$ to match the scheduling of \cite{edm}, as it allows us to take advantage of their optimized choice of time steps and leads to better comparability. Lastly, we introduce a parameter $p>0$ to regulate the relative speed of convergence of $\sigma$ and $\gamma$ via $\gamma(t)=\sigma(t)^p$.

With these choices, (\ref{eq:modified_inference}) simplifies to 
\begin{equation}\label{eq:modified_inference2}
    x_{i+1} = x_{i}- m(t_i)(t_{i+1}-t_i)t_i \big((1+\gamma(t_i))S_\theta(t_i,x_i)+\tfrac{b\gamma(t_i)}{2t_i}\ddsigma S_\theta(t_i,x_i)\big),
\end{equation}
where $\ddsigma S_\theta$ denotes the derivative of $S_\theta\colon\ [0,\sigma_{\rm max}]\times \R^d\to\R^d$ w.r.t\@ its first argument, and $\gamma(t_i)$ is a solution of
\begin{align*}
    \sqrt{a\gamma(t_i)^{\frac{2}{p}}+b\gamma(t_i)}=t_i.
\end{align*}
We observe that in the case of $S_\theta(\sigma,x)=-\tfrac{x-\mu}{\sigma^2}$, i.e.\@ an initial distribution consisting of a Dirac delta at $\mu$,
the standard score inference rule (\ref{eq:basic_Euler}) with $\sigma(t)=t$ evaluates to 
\begin{equation}
\begin{split}
    \label{eq:standard_score_Dirac}
    x_{i+1}&=x_i - \tfrac{t_i-t_{i+1}}{t_i}(x_i-\mu)\\&=(1-\tfrac{t_i-t_{i+1}}{t_i})x_i + \tfrac{t_i-t_{i+1}}{t_i}\mu .
    \end{split}
\end{equation}
This means that, if $\tfrac{t_i-t_{i+1}}{t_i}\ge \beta>0$ for every $i$, we have
\begin{align*}
    \|x_{i+k} -\mu\|\leq(1-\beta)^k\|x_i-\mu\|,
\end{align*}
i.e.\@ convergence to $\mu$ at a geometric rate. We would like our inference rule to mimic this behavior for simple Dirac deltas, which is accomplished by choosing the correction factor as
\begin{align}\label{eq:correction_factor}
    m(t_i)
=\left(1+\gamma(t_i)-\tfrac{b\gamma(t_i)}{t_i^2}\right)^{-1},
\end{align}
see Section~\ref{sec:proofs} in the appendix for a detailed derivation.
Note that, as long as $\gamma(t_i)\in o(t_i^{-2})$ we have $\lim_{t\to 0}m(t)=1$. Putting it all together, we arrive at the following algorithm.
\begin{algorithm*}
    \caption{Inference with extended score} \label{alg:main} 
    \begin{algorithmic}[1]
       \STATE \textbf{function} \textsc{Inference}$(S_\theta,a,b,p,\Delta,(t_i)_{i\in\{0,\dots,N\}})$ 
        \STATE \textbf{sample} $x_0\sim\cN(0,t_0^2\mathbb{I})$ \COMMENT{Generate Gaussian sample}
        \FOR{$i\in\{0,\dots,N-1\}$}
            \STATE  $\gamma_i =  \mathrm{\textbf{solve}}(a\gamma_i^{2/p} + b\gamma_i -t_i^2 =0)$ \COMMENT{Determine $\gamma_i$ using a root finding algorithm}
            \STATE $s_i = S_\theta(t_i,x_i)$ \COMMENT{Evaluate the score network}
            \STATE $\widetilde{s}_i = S_\theta((1+\Delta)t_i,x_i)$
            \STATE $s'_i = \tfrac{\widetilde{s}_i-s_i}{\Delta \cdot t_i}$ \COMMENT{Compute an approximation to the derivative}
            \STATE $m_i=(1+\gamma_i-\tfrac{b\gamma_i}{t_i^2})^{-1}$ \COMMENT{Compute correction factor}
            \STATE $x_{i+1}= x_i-m_i(t_{i+1}-t_i)t_i((1+\gamma_i)s_i+\tfrac{b\gamma_i}{2t_i}s'_i)$ \COMMENT{Update sample}
        \ENDFOR 
        \STATE \textbf{return} $x_N$
    \end{algorithmic}
\end{algorithm*}


% You may use color figures.
% However, it is best for the
% figure captions and the paper body to make sense if the paper is printed
% either in black/white or in color.
% \begin{figure}[h]
% \begin{center}
% %\framebox[4.0in]{$\;$}
% \fbox{\rule[-.5cm]{0cm}{4cm} \rule[-.5cm]{4cm}{0cm}}
% \end{center}
% \caption{Sample figure caption.}
% \end{figure}

% \begin{table}[t]
% \caption{Sample table title}
% \label{sample-table}
% \begin{center}
% \begin{tabular}{ll}
% \multicolumn{1}{c}{\bf PART}  &\multicolumn{1}{c}{\bf DESCRIPTION}
% \\ \hline \\
% Dendrite         &Input terminal \\
% Axon             &Output terminal \\
% Soma             &Cell body (contains cell nucleus) \\
% \end{tabular}
% \end{center}
% \end{table}
%\newpage
\section{Numerical Experiments}\label{sec:numerics}

In section~\ref{sub:51} (toy examples in $\R^2$) and in section~\ref{sub:ffhq} (FFHQ, AFHQv2, and ImageNet) we present some illustrative numerical simulations to build intuition of the ``manifold attraction'' property of the extended score on clean datasets. Further, in sections~\ref{sub:synth} (synthetic dataset) and \ref{sub:cifar} (CIFAR-10) we provide controlled experiments showing qualitatively and quantitatively (FID scores) the denoising effect of the extended score. Finally, in section~\ref{sub:real} we test MAD with real Cryo-EM data. 

We emphasize that MAD addresses a specific blind generative denoising task where only noisy data is available, and the specific degradation model may be
unknown. In this context, there are no established benchmarks, and a comparison with other (non-generative) manifold denoising approaches would not be meaningful, because they tackle a different problem.



\subsection{Illustrative Examples in \texorpdfstring{$\R^2$}{R2}}
\label{sub:51}
\begin{figure}
    \centering
    \includegraphics[width=0.95\linewidth]{trajectories2a.png}
    \includegraphics[width=0.95\linewidth]{trajectories2b.png}
    \caption{Samples (blue) of $1$-dimensional Gaussian mixture supported on an affine subspace with variances $2$, $0.5$, $0.25$, $0.1$, and $0.05$ respectively. Inference trajectories for standard score (top) and extended score (bottom) are indicated in red.} 
    \label{fig:1Dtraj}
\end{figure}


\begin{figure}
    \centering
    \begin{subfigure}{0.5\textwidth}
    \includegraphics[trim={0.5cm 0 0.25cm 0.2cm},clip,width=\linewidth]{tiltedGaussians.png}
    \subcaption{Tilted Gaussians}\label{fig:2DGaussians_a}
    \end{subfigure}
    \begin{subfigure}{0.5\textwidth}
    \includegraphics[trim={0.5cm 0 0.25cm 0.2cm},clip,width=\linewidth]{radialGaussian2.png}      
    \subcaption{Radial Gaussians}\label{fig:2DGaussians_b}
    \end{subfigure}
    \caption{Target distribution is displayed as a histogram colormap. Red crosses indicate samples generated via standard (left) or extended (right) score inference. Extended score parameters are $p=1.3$, $a=1$, $b=1.1$ for (a) and $p=2$, $a=2$, $b=30$ for (b). 
    %Comparison of inference results for Gaussian mixture of the form $\sum_{i=1}^N \tfrac{1}{2\pi N} \det(\Sigma_i)^{-\frac{1}{2}}e^{-\frac{1}{2}(x-\mu_i)^T\Sigma_i^{-1}(x-\mu_i)}$ with $\Sigma_i=R_i\begin{pmatrix}1.5 &0\\0 & 0.05 \end{pmatrix}$, where the $R_i$ are randomly chosen rotation matrices. Extended score is applied with parameters ... .
    }
    \label{fig:2DGaussians}
\end{figure}

 We visualize the evolution according to the extended score in $\R^2$ for relatively simple distributions, where the (extended) score can be computed explicitly (see Appendix~\ref{sec:product}). Figure~\ref{fig:1Dtraj} displays the inference trajectories for a degenerate Gaussian mixture, which is supported on an $1$-dimensional affine subspace, i.e.\@ it behaves like a Dirac in the off-manifold direction. The means are chosen such that the different mixture components affect each other very little.
 %behaves like a Dirac delta at $0$ in the $x_2$-direction and a Gaussian in the $x_1$-direction as the means are chosen such that the different mixture components affect each other very little.
 We can see that for the Gaussians with low variance using the extended score causes all trajectories to end up at the respective means, i.e.\@ it behaves as if we had Dirac deltas (in the on-manifold direction) at these locations. For the higher variance Gaussians the extended score trajectories still end up closer to the mean than they would for standard score, but this variance reduction effect decreases significantly as the variance of the Gaussians increases. All together, we essentially have a soft thresholding effect, where variances below a certain value are shrunk to $0$, while large enough variances are left almost unchanged.    

  
% \begin{figure*}
%     \centering
%     \includegraphics[width=0.95\linewidth]{1Dtraj.png}
%     \caption{Comparison of inference trajectories (red) and samples (blue) for Gaussian mixture {\small $p(x_1,x_2)=\delta(x_2)\sum_{i=1}^5 \tfrac{1}{5}(2\pi v_i)^{-\frac{1}2{}}\mathrm{exp}(-\frac{(x_1-\mu_i)^2}{2v_i})$} with equal weights, variances $v=(0.2,0.5,1,2,4)$, and means $\mu=(-20,-10,0,10,20)$. Extended score is applied with parameters $a=b=p=1$.} 
%     \label{fig:1Dtraj}
% \end{figure*}




In Figure~\ref{fig:2DGaussians_a} we consider a mixture of Gaussians with covariance matrices  $\Sigma_i=R_i\mathrm{diag}(1.7,0.2)$, where the $R_i$, $i\in\{1,\dots,21\}$, are randomly chosen rotation matrices, i.e.\@ around each mean of the mixture we have large variance in one direction and small variance in the direction orthogonal to the first one. Locally this can be viewed as a  $1.5$-variance Gaussian along some $1$-dimensional affine linear subspace to which Gaussian noise with covariance matrix $0.2\mathbb{I}$ is added, resulting in a $1.7$ variance in one direction and a $0.2$ variance in the other. The denoising effect is clear: the extended score inference moves the points onto the affine linear subspace corresponding to the first principal direction of each Gaussian, while leaving the spread along the affine linear subspace almost unchanged.
In Figure~\ref{fig:2DGaussians_b} we demonstrate that this effect is not limited to affine linear manifolds, by considering a mixture of radial Gaussians of the form {\small $p(x) = C\sum_{i=1}^5 \mathrm{exp}(-\tfrac{(\|x-\mu_i\|-r)^2}{2v})$} with variance $v=2.5$, radius $r=10$, randomly chosen means $\mu_i$, and $C$ a normalizing constant.


\subsection{Effects on Noiseless FFHQ, AFHQv2, and ImageNet}

\label{sub:ffhq}



While the last section provides some basic intuition for the effects of our inference method, it is, of course, unclear how this translates to much more complicated distributions in much higher dimensions. We will explore this question using our inference method with the pretrained score networks from \cite{edm} for\footnote{Note that these networks have been trained on images from these datasets which have been downsampled to  $64\!\times\! 64$ resolution.} FFHQ (\cite{FFHQ}), AFHQv2 (\cite{AFHQ}), and ImageNet (\cite{ImageNet}), where we also use the time schedule suggested in \cite{edm}, and only vary the hyperparameters specific to our method, i.e.\@ $a$, $b$, $p$, and $\Delta$ in Algorithm~\ref{alg:main}. As these datasets are arguably essentially noiseless, the effect of the extended score will be visible on the different features of the images. The following sections consider noisy data.

\begin{figure}[ht]
    \centering
    \includegraphics[trim={0 6.775cm 0 0},clip,width=0.98\linewidth]{FFHQ_main.png}
    
    \includegraphics[trim={0 6.775cm 0 0},clip,width=0.98\linewidth]{AFHQ_main.png} 
    \caption{In both subfigures, all images in each row start from the same latent noise sample, and the leftmost column uses standard score whereas second to last columns use Algorithm~\ref{alg:main} with $\Delta=0.0001$, $p=8$, $a=2.5$, and, left to right, $b\in\{2,5,10,20,40,80\}$. Further examples in Figure~\ref{fig:faces_appendix}.}\label{fig:faces_main}
\end{figure}

In Figure~\ref{fig:faces_main}, we showcase the impact of our method on the generation of human or animal faces using the pretrained score networks for FFHQ and AFHQv2, respectively. Note that increasing the parameter $b$ in Algorithm~\ref{alg:main} results in a greater impact of the extended score, i.e.\@ we expect stronger attraction to the manifold of primary variation. We observe that for all values of $b$, we generate samples with qualitatively the same facial features, but for larger values of $b$, we always get a plain single-colored background. 
In case of faces, it seems quite clear that the direction of primary relevance should correspond to essential facial features that are present in all the data, whereas the background variation, as well as features like glasses and head wear, are split across a much larger number of directions and are thresholded out first. Furthermore, as illustrated in the appendix (see Figures~\ref{fig:FFHQ_ab} and \ref{fig:faces_appendix}), the method demonstrates improved stability across hyperparameter variations on FFHQ compared to simpler synthetic datasets. This indicates that the higher complexity of natural images, reflected into the higher dimensionality of the underlying manifold, may inherently improve hyperparameter robustness.
We  also observe similar effects on ImageNet, e.g.\@ in the examples in Figure~\ref{fig:ImageNetmain}, where the extended score inference seems to focus on generating one primary object, while progressively thresholding out everything else as we increase $b$. However, since ImageNet contains a large variety of objects, the primary directions differ between classes, and thus the effect of extended score inference with a given choice of hyperparameters is much more varied across different starting noise images (see Figure~\ref{fig:INET_appendix} in the Appendix).  

\begin{figure}
    \centering
    \includegraphics[width=0.99\linewidth]{INET_horse.png}
    \includegraphics[width=0.99\linewidth]{INET_mushroom.png}
    \caption{Each row starts from the same latent noise sample, and the leftmost column uses standard score whereas second to last columns use Algorithm~\ref{alg:main} with $\Delta=0.001$, $p=12$, and $a=8$, $b\in\{5,20,50,80,250\}$ for the top row, $a=4$, $b\in\{2,10,30,60,100\}$ for the bottom row.}\label{fig:ImageNetmain}
\end{figure}


Although Algorithm~\ref{alg:main} is deterministic when started with the same random seed, i.e.\@ such that it starts from the same latent noise image, it relies on a score network trained through a highly stochastic process. Consequently, even small changes of the iterates $x_i$ can build up and lead to the generation of a significantly different image. As can be seen in Figure~\ref{fig:faces_main}, in general the algorithm is not particularly suited for orthogonal projection on the manifold of, in this case, faces, i.e.\@ it does not simply generate the same face as standard inference would but without background. However, as showcased in Figure~\ref{fig:FFHQ_ab} in the Appendix, there are often certain ranges of parameters that lead to rather similar images and may be used for manual adjustment of a generated image. 



\begin{figure}
    \centering
    \includegraphics[trim={0 2.2cm 0 0},clip,width=0.98\linewidth]{synth_clean2.png}
    
    \includegraphics[trim={0 2.2cm 0 0},clip,width=0.98\linewidth]{synth_noisy2.png}
    \caption{Grid of clean samples from the synthetic data set (top); grid of corresponding noisy samples (bottom). Each sample is an $8$-bit grayscale $64\!\times\!64$ image, displayed via viridis colormap.}
    \label{fig:synth_data}
\end{figure}


\subsection{Denoising Synthetic Data}

\label{sub:synth}
% \begin{figure}
%     \centering
%     \includegraphics[width=0.43\linewidth]{synth_clean.png}
%     \includegraphics[width=0.43\linewidth]{synth_noisy.png}
%     \includegraphics[width=0.12\linewidth]{colourbar.png}
%     \caption{Grid of clean samples from the synthetic data set (left); grid of corresponding noisy samples (right). Each sample is an $8$-bit grayscale $64\!\times\!64$ image, displayed via viridis colormap.}
%     \label{fig:synth_data}
% \end{figure}


 
\begin{figure}
    \centering
    \includegraphics[trim={0 2.2cm 0 0},clip,width=0.98\linewidth]{synth_extended_joint.png}
    
    \includegraphics[trim={0 2.2cm 0 0},clip,width=0.98\linewidth]{synth_standard2.png}
    \caption{Grid of samples generated from consecutive random seeds by Algorithm~\ref{alg:main} with parameters $\Delta=0.02$, $p=8$, $a=0.002$, and $b=15$ (top); grid of samples generated by inference with standard score (bottom). Each sample is an $8$-bit grayscale $64\!\times\!64$ image, displayed via viridis colormap.}
    \label{fig:synth_results}
\end{figure}


We created a synthetic dataset of clean samples, added noise, and trained a diffusion model only on the noisy data. Our clean distribution consists of $8$-bit grayscale $64\!\times\!64$ images depicting $4$ different shapes, each appearing with equal probability and rotated by an angle chosen uniformly at random (see Figure~\ref{fig:synth_data}, top). As such, this probability distribution is essentially %\footnote{Due to discretization effects its support is, strictly speaking, not a union of four $1$-dimensional manifolds in pixel space, but very closely concentrated around one.\de{**not sure how necessary/helpful this comment is**} \ga{** I'd be happy to remove it, I don't think it's relevant for our aims **}} 
supported on $4$ disjoint $1$-dimensional manifolds in pixel space. We then add two types of corruption to obtain noisy samples. First, we blot out large parts of the shape, by uniformly randomly picking $50$ locations on the boundary of the shape and subtracting Gaussian bump functions centered at those locations from the clean image. Second, we add i.i.d.\@ Gaussian noise to each pixel, resulting in images where the original shapes can hardly be discerned from a single image (see Figure~\ref{fig:synth_data}, bottom). We then trained a DDPM++ model\footnote{Using the method from the github to \cite{edm} with default settings except for $\mathrm{cond}=0$ and $\mathrm{augment}=0$, which took roughly $14$ hours on a single $A100$ GPU.} on a set of $100\,000$ of  noisy samples for a duration of $3$~mimg, i.e.\@ $30$ repeats per image. When using the trained network for standard score inference we obtain a reproduction of the corrupted samples (see Figure~\ref{fig:synth_results}, bottom). In contrast, running our extended score inference, we obtain samples showing the original shapes (see Figure~\ref{fig:synth_results}, top). Figure~\ref{fig:synth_ab_grid} in the Appendix illustrates the dependence on the extended score parameters $\Delta$, $a$, $b$, and $p$. In order to test whether our method works on corrupted training data beyond the case of Gaussian noise, we also trained a DDPM++ model, using the same training settings, on a set of $100\,000$ synthetic samples, where each pixel was set to $0$ with probability $0.5$ (see Figure~\ref{fig:synth-pixel} in Appendix \ref{sub:additionalnumerics}). As can be seen in Table~\ref{table}, for both cases extended score inference leads to a significant improvement in FID compared to the standard score, which generates samples with essentially the same distance to the clean data as the noisy data.

\begin{table*}
\begin{center}
\begin{tabular}{ |l|c|c|c| } 
 \hline
Noisy dataset  & training data & standard generation & extended generation\\
 \hline
 synthetic data, Gaussian noise & $320.56$ & $319.42$ &  $\textbf{189.51}$\\ 
 synthetic data, pixel removal & $233.86$ & $234.40$ &  $\textbf{163.14}$\\ 
 CIFAR-$10$, Gaussian noise & $136.87$ & $147.47$ &  $\textbf{77.15}$\\ 
 FFHQ, Gaussian noise & $133.60$ & $134.42$ &  $\textbf{74.52}$\\ 
 FFHQ, Shot noise & $257.28$ & $257.42$ &  $\textbf{101.12}$\\ 
 FFHQ, pixel removal & $372.21$ & $387.87$ &  $\textbf{104.00}$\\ 
 \hline
\end{tabular}
\end{center}
\caption{FIDs with respect to the corresponding clean dataset, computed from $10\,000$ images.
\label{table}}
\end{table*}



 
% \begin{figure}
%     \centering
%     \includegraphics[width=0.43\linewidth]{synth_extended.png}
%     \includegraphics[width=0.43\linewidth]{synth_standard.png}
%     \includegraphics[width=0.12\linewidth]{colourbar.png}
%     \caption{Grid of samples generated by Algorithm~\ref{alg:main} with parameters $\delta=0.02$, $p=8$, $a=0.002$, and $b=15$ (left); grid of samples generated by inference with standard score (right). Each sample is an $8$-bit grayscale $64\!\times\!64$ image, displayed via viridis colormap. \ga{** Perhaps we've already discussed this, but the angles of the generated images are almost always the same, and not uniform in $0-2\pi$. Is there a way to slightly modify the parameters and get more uniform angles? **}}
%     \label{fig:synth_results}
% \end{figure}

\subsection{Denoising CIFAR-10 and FFHQ}
\label{sub:cifar}
\begin{figure}
    \centering
    \includegraphics[trim={0 7.91cm 0 0},clip,width=0.98\linewidth]{rev_cifar_modified.png}
    
    \includegraphics[trim={0 7.91cm 0 0},clip,width=0.98\linewidth]{rev_cifar_euler.png}
    \caption{Standard score inference (bottom) and extended score inference (top) for CIFAR-$10$ with additive Gaussian noise.}
    %Trained on CIFAR-$10$ with additive Gaussian noise and generated by standard score inference (bottom) or by extended score inference (top).}
    \label{fig:cifar_additive_Gaussian}
\end{figure}


Next, we test whether the denoising capabilities of our method extend to real data. Specifically we train DDPM++ models on CIFAR-$10$ and FFHQ with additive Gaussian noise. While standard score inference produces image with roughly the amount of noise of the images it was trained on, extended score inference produces significantly cleaner images, both visually (Figure~\ref{fig:cifar_additive_Gaussian} and Figure~\ref{fig:ffhq_additive_Gaussian}) as well as quantitatively in the form of FID scores (Table~\ref{table}). 
We note that the FID values for standard generation very closely match the values for the respective training sets, indicating that in training the diffusion model simply learns to reproduce the corruption and does contribute to the denoising effect.
The intermediate images generated during inference (see Figure~\ref{fig:paths} in Appendix~\ref{sub:additionalnumerics}) strongly indicate that this effect cannot simply be achieved by truncated sampling. 

We also train DDPM++ models on FFHQ corrupted by shot/Poisson noise (Figure~\ref{fig:ffhq_PN}) or pixel removal noise (Figure~\ref{fig:ffhq_PR}).
While the images generated by extended score inference are of somewhat worse quality than in the Gaussian case, the FID values of the training sets are significantly larger, i.e.\@ we are dealing with a higher degree of corruption. In fact, MAD achieves a larger FID improvement over standard generation than in the Gaussian case.
%the relative improvement in FID compared to training set resp.\@ standard generation is significantly higher (see Table~\ref{table}).
More importantly, it demonstrates that MAD is effective for non-Gaussian corruptions, which is evidence that it does not simply modify the amount of Gaussian noise removed during inference.


\begin{figure}
    \centering
    \includegraphics[trim={0 15.8cm 0 0},clip,width=0.98\linewidth]{ffhq_GN_modified.png}
    
    \includegraphics[trim={0 15.8cm 0 0},clip,width=0.98\linewidth]{ffhq_GN_euler.png}
    \caption{Standard score inference (bottom) and extended score inference (top) for FFHQ with additive Gaussian noise.}
    \label{fig:ffhq_additive_Gaussian}
\end{figure}

\begin{figure}
    \centering
    \includegraphics[trim={0 15.8cm 0 0},clip,width=0.98\linewidth]{ffhq_PN_modified.png}
    
    \includegraphics[trim={0 15.8cm 0 0},clip,width=0.98\linewidth]{ffhq_PN_euler.png}
    \caption{Standard score inference (bottom) and extended score inference (top) for FFHQ with Shot/Poisson noise.}
    \label{fig:ffhq_PN}
\end{figure}



\begin{figure}
    \centering
    \includegraphics[trim={0 15.8cm 0 0},clip,width=0.98\linewidth]{ffhq_PR_modified.png}
    
    \includegraphics[trim={0 15.8cm 0 0},clip,width=0.98\linewidth]{ffhq_PR_euler.png}
    \caption{Standard score inference (bottom) and extended score inference (top) for FFHQ with pixel removal noise.}
    \label{fig:ffhq_PR}
\end{figure}



\subsection{Denoising Real Data}
\label{sub:real}
In this section we test our method on real data from single-particle Cryo-Electron Microscopy (see \cite{CryoEMprimer} for an introduction), where many particles of the same type are suspended in liquid, frozen, and put under an electron microscope. This produces extremely noisy $2$D-images which need to be refined before further steps like $3$D reconstruction can be attempted.
This presents an opportunity to investigate the performance of our method on practically relevant real world data with non-Gaussian noise. We use the EMPIAR-$11618$ \citep{EMPIARpaper} dataset of $68\,401$ grayscale images with  $256\!\times\!256$ resolution. They were extracted from  raw data and undergone some preprocessing, but are still very noisy (see Figure~\ref{fig:real_dataset}, top). We trained a DDPM++ model\footnote{In order to compensate for the higher resolution we reduced the number of feature channels in the DDPM++ architecture from $128$ to $32$, but otherwise used the same settings as for the synthetic data. Training took roughly $60$ hours on two $A100$ GPUs.} on this data and used our method to generate samples (see Figure~\ref{fig:real_dataset}, middle) whose shapes correspond strongly to what has been obtained by \cite{EMPIARpaper}, see Figure~\ref{fig:real_dataset}, bottom. We emphasize that the network has only ever seen noisy data and has in no way been specifically adjusted based on a priori knowledge of these shapes. This can be seen by the fact that standard score inference generates noisy samples, very similar to those in the training set (see Figure~\ref{fig:real_generated} in the Appendix). The parameters used for Figure~\ref{fig:real_dataset} of the extended score inference have been determined by hand with such knowledge, of course, but similar results are generated for a wide range of parameter choices (see Figure~\ref{fig:real_ab_grid} in the Appendix). In particular, this demonstrates a significant capability of the extended score to guide generation towards samples from the underlying image manifold, also in the case of real data with extreme noise corruption caused by physical measurement modalities. This example serves as proof of concept for the validation of our approach: reaching state-of-the-art performance is outside the scope of this work, and would require incorporating more prior domain knowledge.


% \begin{figure}
%     \centering
%     \begin{subfigure}{0.432\textwidth}
%         \includegraphics[width=\linewidth]{real_data36.png}     
%     \end{subfigure}
%     \begin{subfigure}{0.432\textwidth}
%         \includegraphics[width=\linewidth]{real_extended.png}     
%     \end{subfigure}
%     \includegraphics[width=0.12\linewidth]{colourbar.png}\\
%     \includegraphics[width=0.161\linewidth]{emd1.jpeg}
%     \includegraphics[width=0.161\linewidth]{emd2.jpeg}
%     \includegraphics[width=0.161\linewidth]{emd3.jpeg}
%     \includegraphics[width=0.161\linewidth]{emd4.jpeg}
%     \includegraphics[width=0.161\linewidth]{emd5.jpeg}
%     \includegraphics[width=0.161\linewidth]{emd6.jpeg}
%     \caption{Images from the EMPIAR-$11618$ dataset (left). Samples generated by Algorithm~\ref{alg:main} with parameters $\delta=0.001$, $p=8$, $a=0.01$, and $b=2$ (right). Shapes obtained by \cite{EMPIARpaper}(bottom). The colormap is applied with normalization per image to enhance contrast.}
%     \label{fig:real_dataset}
% \end{figure}

\begin{figure}
    \centering
    \begin{subfigure}{0.492\textwidth}
        \includegraphics[trim={0 6.77cm 0 0},clip,width=\linewidth]{real_data40.jpeg}     
        %\includegraphics[trim={0 9cm 0 0},clip,width=\linewidth]{real_data40.png}     
    \end{subfigure}
    \begin{subfigure}{0.492\textwidth}
        \includegraphics[trim={0 6.77cm 0 0},clip,width=\linewidth]{real_extended2.jpeg}   
        %\includegraphics[trim={0 9cm 0 0},clip,width=\linewidth]{real_extended2.png}     
    \end{subfigure}\\
    \includegraphics[width=0.15\linewidth]{emd1.jpeg}
    \includegraphics[width=0.15\linewidth]{emd2.jpeg}
    \includegraphics[width=0.15\linewidth]{emd3.jpeg}
    \includegraphics[width=0.15\linewidth]{emd4.jpeg}
    \includegraphics[width=0.15\linewidth]{emd5.jpeg}
    \includegraphics[width=0.15\linewidth]{emd6.jpeg}
        \caption{Images from the EMPIAR-$11618$ dataset (top). Samples generated by Algorithm~\ref{alg:main} with  $\Delta=0.001$, $p=8$, $a=0.01$, $b=2$ (middle). Shapes obtained by \cite{EMPIARpaper} (bottom), taken from \citep[EMD-17944]{EMDB}. The colormap is applied with normalization per image to enhance contrast.}
    \label{fig:real_dataset}
\end{figure}


%\newpage
\section{Conclusion}

%In this work, we introduced Manifold Attracted Diffusion (MAD), a novel modification to the inference procedure of score-based diffusion models that enables approximate sampling from clean data distributions despite training on noisy datasets.
We introduced Manifold Attracted Diffusion (MAD), a novel inference approach for score-based diffusion models to generate clean samples from a distribution despite training on noisy datasets.
%By leveraging the manifold hypothesis and defining an extended score, our method suppresses off-manifold noise while preserving on-manifold variations, which results in attracting samples toward a low-dimensional structure.
Based on the manifold hypothesis, our method utilizes the underlying geometry to suppress off-manifold variations while preserving on-manifold variations, which results in attracting samples toward a low-dimensional structure.
%It is an inference-only modification that can be applied to many pre-trained diffusion models without requiring retraining. The computational overhead is minimal, adding only one extra evaluation of the score network per step to approximate a derivative.
% It is an inference-only modification that can be applied to many pre-trained diffusion models without requiring retraining. The computational overhead is minimal, adding only one extra evaluation of the score network per step to approximate a derivative.
It can leverage established training algorithms and pretrained models.
The required computation cost is at most twice that of standard inference, due to requiring a second evaluation of the score network. However, we observed empirically that it is enough to use the extended score only in the last half of the inference steps, yielding an additional computational cost of 50\%, but additional investigation on this aspect is required.
% Our experiments on both toy examples and high-dimensional image data demonstrate that MAD successfully reduces variance and suppresses noise.
Numerical experiments on both synthetic and real data demonstrate that MAD successfully suppresses noise.



Future directions include extending MAD to solve inverse problems, as in Cryo-EM for denoising individual images, or for other denoising or image restoration tasks.
This would benefit from integration with conditional diffusion models %, as explored in diffusion posterior sampling frameworks 
\citep{chungdiffusion, chung2022improving}. Developing automatic and adaptive parameter selection for $\gamma(t)$ would enhance the methods robustness and applicability. Further theoretical analysis of the extended score may lead to an improved incorporation into the inference procedure, in particular one could combine it with inference techniques using noise injection or higher order ODE discretization schemes. Future work should also involve a rigorous quantitative study comparing the performance and computational trade-offs of MAD against training-time methods, finetuning on limited clean data, and other relevant inference-time baselines such as truncated sampling \citep{daras2025how}. Furthermore, it would be valuable to investigate whether the MAD framework could be adapted to time-unconditioned generative models \citep{sun2025is}, or if our method fundamentally relies on a time-dependent score. It may also be interesting to consider whether the MAD framework can be adapted for Variance Preserving diffusion models. The key difficulty would be that one cannot use identity \eqref{eq:VE_trick} in that setting. Finally, a promising future direction, motivated by the conceptual similarities to classifier-free guidance \citep{ho2021classifierfree}, is to explore whether our extended score can serve as a more interpretable, geometrically-grounded method for improving generation quality, even for models trained on clean data.


\section*{Acknowledgments}
We would like to thank Pavol Harar for his advice on the CryoEM example.
This material is based upon work supported by the Air Force Office of Scientific Research under award numbers FA8655-23-1-7083. Co-funded by the European Union (ERC, SAMPDE, 101041040  and MSCA, CompAS, 101152591). Views and opinions
expressed are however those of the authors only and do not necessarily reflect those of the European
Union or the European Research Council. Neither the European Union nor the granting authority
can be held responsible for them. 
The research was supported in part by the MIUR Excellence Department Project awarded to Dipartimento di Matematica, Università di Genova, CUP D33C23001110001. Co-funded by European Union – Next Generation EU, Missione 4 Componente 1 CUP D53D23005770006. The numerical results have been achieved with the help of the Austrian Scientific Computing (ASC) infrastructure.

\section*{Conflict of Interest}
The authors declare no conflict of interest.

\section*{Impact Statement}
This paper presents work whose goal is to advance the field of Machine Learning. There are many potential societal consequences of our work, none which we feel must be specifically highlighted here.

\bibliography{iclr2026_conference}
\bibliographystyle{icml2026}

\newpage
\appendix
\onecolumn


\section{The Extended Score for Products of Measures}\label{sec:product}

In the following lemma, we derive the expression of the extended score of the product of two measures. In this section, the superscripts denote the dimension of the domain.

\begin{lemma}\label{lem:score-product}
    Let $d_1+d_2=d$. Take $p_1\in \widetilde{M}(\R^{d_1})$  and $p_2\in \widetilde{M}(\R^{d_2})$. Then the product measure $p=p_1\otimes p_2$ belongs to $\widetilde{M}(\R^{d})$ and
    \[
H^d_0 p(x)=(H^{d_1}_0 p_1(x_1),H_0^{d_2} p_2(x_2)),
    \]
    where $x=(x_1,x_2)\in \R^{d_1+d_2}$.
\end{lemma}
\begin{proof}
  We have that $g_\gamma^d (x) = g_\gamma^{d_1}(x_1)g_\gamma^{d_2}(x_2)$. Thus, by Fubini's theorem, we have
  \[
(p*g^d_\gamma)(x) = (p_1*g^{d_1}_\gamma)(x_1)\cdot (p_2*g^{d_2}_\gamma)(x_2).
  \]
  Thus
    \[
\log (p*g^d_\gamma)(x) =\log  (p_1*g^{d_1}_\gamma)(x_1)+\log (p_2*g^{d_2}_\gamma)(x_2).
  \]
  Taking a gradient with respect to $x$ we obtain
  \[
  S^d(p*g^d_\gamma)(x) = (S^{d_1}(p_1*g^{d_1}_\gamma)(x_1),0_{d_2})+(0_{d_1},S^{d_2}(p_2*g^{d_2}_\gamma)(x_2)).
  \]
  Therefore, since the expression of $H^d_\gamma (p)$ is linear in $S^d(p*g^d_\gamma)$, we obtain
\[
H^d_\gamma (p)(x) =( H^{d_1}_\gamma (p_1)(x_1),0_{d_2})+(0_{d_1},H^{d_2}_\gamma (p_2)(x_2))=( H^{d_1}_\gamma (p_1)(x_1),H^{d_2}_\gamma (p_2)(x_2)).
\]
Taking the limit as $\gamma\to 0$, the result follows.
\end{proof}


This result can be used to calculate the extended score of a degenerate distribution $p$, namely, a distribution supported on a lower-dimensional affine subspace of $\R^d$. Since the score, and thus the extended score, is equivariant with respect to rotations and translations, without loss of generality, we can assume that $p$ is supported on $\{(x_1,0_{d_2})\in\R^d:x_1\in\R^{d_1}\}$ with density $p_1\in P(\R^{d_1})$, namely,
\[
p=p_1\otimes \delta_{d_2},
\]
where $\delta_{d_2}\in\widetilde{M}(\R^{d_2})$ is the Dirac delta centered at $0$ in $\R^{d_2}$. By Lemmata~\ref{lem:extension} and \ref{lem:score-product}, we obtain that the extended score of $p$ is given by
\[
H_0^d p(x) = (H^{d_1}_0 p_1(x_1),H_0^{d_2} p_2(x_2))
= (S^{d_1} p_1(x_1),-x_2),
\]
where we have also used that $H_0\delta(x)=-x$.

In the particular case when $p_1$ is a non-degenerate Gaussian distribution on $\R^{d_1}$ with mean $\mu_1\in\R^{d_1}$ and covariance $\Sigma_1\in\R^{d_1\times d_1}$, namely $p_1=\mathcal{N}(\mu_1,\Sigma_1)$, we have
\[
H_0^d p(x) = (-\Sigma_1^{-1}(x_1-\mu_1),-x_2).
\]
As expected, this coincides with the standard score of the (non-degenerate) Gaussian distribution $\mathcal{N}((\mu_1,0_{d_2}),\Sigma)$ on $\R^d$, where $\Sigma$ is the block matrix given by
\[
\Sigma= \begin{bmatrix}
    \Sigma_1 & 0\\
    0& \mathbb{I}^{d_2}
\end{bmatrix}.
\]



\section{Proofs and Derivations}\label{sec:proofs}
\begin{proof}[Proof of Lemma~\Ref{lem:extension}]
    By standard properties of approximate identities \citep[{Example~1.2.17 and Theorem~1.2.19(2)}]{grafakos-2014},
    for $f\in C(\R^d)$ bounded we have that
    \begin{align*}
        \lim_{\gamma\to 0} (f*g_\gamma)(x)
        = f(x),\qquad x\in\R^d,
    \end{align*} 
    Analogously, with $h_\gamma(x):=\tfrac{1}{\gamma}\|x\|^2 g_\gamma(x)$ by  \citep[Theorem~1.2.21(b)]{grafakos-2014}  we get that
    \begin{align*}
        \lim_{\gamma\to 0} (f*h_\gamma)(x) = C f(x),\qquad x\in\R^d,
    \end{align*}
    where $C=\int_{\R^d} (2\pi)^{-\frac{d}{2}} \|z\|^2e^{-\frac{\|z\|^2}{2}}  \d z < \infty$ . In addition, we observe that
    \begin{align*}
        \ddxi g_\gamma(x)=-\tfrac{x_i}{\gamma}g_\gamma(x)
    \end{align*}
    and
    \begin{align*}
        \ddgamma g_\gamma(x) = (\tfrac{-d}{2\gamma}+\tfrac{\|x\|^2}{2\gamma^2})g_\gamma(x) = \tfrac{1}{2\gamma}(h_\gamma(x)-dg_\gamma(x)). 
    \end{align*}
    Consequently, as $p(x)>0$ for all $x\in\R^d$, by assumption
    \begin{align*}
        \lim_{\gamma\to 0} \gamma\ddgamma S&(p*g_\gamma)
        = \lim_{\gamma\to 0} \gamma\ddgamma\frac{\nabla_x(p*g_\gamma)}{p*g_\gamma}=\lim_{\gamma\to 0} \gamma\ddgamma\frac{\nabla_x p*g_\gamma}{p*g_\gamma}\\
        &=\lim_{\gamma\to 0} \gamma \frac{(p*g_\gamma)(\nabla_xp*\ddgamma g_\gamma)-(\nabla_x p*g_\gamma)(p*\ddgamma g_\gamma)}{(p*g_\gamma)^2}\\
        &=\lim_{\gamma\to 0} \frac{(p*g_\gamma)\big((\nabla_x p* h_\gamma)-d(\nabla_x p* g_\gamma)\big)-(\nabla_x p*g_\gamma)\big(( p* h_\gamma)-d( p* g_\gamma)\big)}{2(p*g_\gamma)^2}\\
        &=\frac{p(C\nabla_x p -d\nabla_x p)- \nabla_x p (Cp-dp)}{2p^2}\\
        &=0.
    \end{align*}
    Moreover, it holds that
    \begin{align*}
        \lim_{\gamma\to 0}S(p*g_\gamma) = \lim_{\gamma\to 0} \frac{\grad{x}(p*g_\gamma)}{p*g_\gamma}=\frac{\lim_{\gamma\to 0}\grad{x}p*g_\gamma}{\lim_{\gamma\to 0}p*g_\gamma}=\frac{\grad{x}p}{p}=Sp,
    \end{align*}
    which completes the proof.
    % \begin{align*}
    %     \ddgamma g_\gamma(x) = (\tfrac{1}{2\gamma} + \tfrac{x^2}{2\gamma^2})g_\gamma(x).
    % \end{align*}
    % We observe the following basic fact about convolution. Let $h_1(x)=x$, $h_2(x)=x^2$, and $f,g\in L^{2}(\R^d)$, then
    % \begin{align*}
    %     (f*)
    % \end{align*}
    % We observe that
    % \begin{align*}
    %     \ddgamma S(p*g_\gamma)=\frac{(\grad{x}p * \ddgamma\g)(p*g_\gamma)-(\grad{x}p *\g)(p*\ddgamma g_\gamma)}{(p*g_\gamma)^2}
    % \end{align*}
\end{proof}
\begin{proof}[Proof of Lemma~\Ref{lem:Dirac_mix}]
     We write $h_i(x):=c_ie^{-\frac{\|x-\mu_i\|^2}{2\gamma}}$, i.e.\@ $p*g_\gamma = \sumi (2\pi\gamma)^{-\frac{d}{2}}  h_i$, 
     and observe that
     \begin{align*}
        S(p*g_\gamma)(x)=-\sum_{i\in[n]}\tfrac{x-\mu_i}{\gamma}w_i(x),
     \end{align*}
     where
     \begin{align*}
         w_i:=\tfrac{h_i}{\sumj{h_j}}.
     \end{align*}
     We observe that
     \begin{align*}
         \ddgamma w_i(x)=\frac{(\sumj h_j(x))\tfrac{\|x-\mu_i\|^2}{2\gamma^2}h_i(x)-h_i(x)\sumj\tfrac{\|x-\mu_j\|^2}{2\gamma^2}h_j(x)}{(\sumj h_j(x))^2}
     \end{align*}
     and
     \begin{align*}
         H_\gamma p(x)
         =(1+\gamma)S(p*g_\gamma)(x) + \gamma \ddgamma S(p*g_\gamma)(x)=\gamma S(p*g_\gamma)(x) +  \ddgamma \gamma S(p*g_\gamma)(x).
    \end{align*}
    
 We will first show that the second term vanishes for $\gamma\to 0$. To this end, we note
    \begin{align*}
         \ddgamma \gamma S(p*g_\gamma)(x)
         &=- \sumi(x-\mu_i)\frac{(\sumj h_j(x))\tfrac{\|x-\mu_i\|^2}{2\gamma^2}h_i(x)-h_i(x)\sumj\tfrac{\|x-\mu_j\|^2}{2\gamma^2}h_j(x)}{(\sumj h_j(x))^2}\\
         &=-\frac{\sum_{i,j\in[n]} h_i(x)h_j(x)(x-\mu_i)(\|x-\mu_i\|^2-\|x-\mu_j\|^2)}{2\gamma^2(\sumj h_j(x))^2}.
     \end{align*}
     Let $x\in W_k$ and $i,j\in[n]$ such that $x\notin W_i \vee x\notin W_j $, then
     \begin{align*}
         \lim_{\gamma\to 0} \frac{h_i(x)h_j(x)}{\gamma^2(\sumj h_j(x))^2}
         \leq \lim_{\gamma\to 0} \frac{h_i(x)h_j(x)}{\gamma^2h_k(x)^2}
         = \lim_{\gamma\to 0} \tfrac{c_ic_j}{\gamma^2 c_k^2}e^{\frac{2\|x-\mu_k\|^2-\|x-\mu_i\|^2 -\|x-\mu_j\|^2}{2\gamma}}=0
     \end{align*}
     as $2\|x-\mu_k\|^2-\|x-\mu_i\|^2 -\|x-\mu_j\|^2 <0$ by definition of $W_i$. Since $\|x-\mu_i\|^2-\|x-\mu_j\|^2 = 0$ if\footnote{In particular, if $i=j$.} $x\in W_i\cap W_j$, and $h_i(x)\geq0$ for every $x\in\R^d$, $i\in[n]$, we have
     \begin{align*}
         \lim_{\gamma\to 0}\ddgamma \gamma S(p*g_\gamma)(x)=0.
     \end{align*}
     
     We proceed by noting that
    \begin{align*}
        \lim_{\gamma\to 0}\frac{h_j(x)}{h_i(x)}
        =\lim_{\gamma\to 0}\tfrac{c_j}{c_i}e^{-\frac{\|x-\mu_j\|^2-\|x-\mu_i\|^2}{2\gamma}}
        =\begin{cases}
            0, & \|x-\mu_i\| < \|x-\mu_j\|\\
            \tfrac{c_j}{c_i} & \|x-\mu_i\| = \|x-\mu_j\|\\
            \infty, & \|x-\mu_i\| > \|x-\mu_j\|
        \end{cases}
    \end{align*}
    and, consequently, using the conventions that $\tfrac{1}{0}=\infty$ and $\tfrac{1}{\infty}=0$,
    \begin{align*}
    %\begin{split}\label{eq:zlim}
        \lim_{\gamma\to 0} w_i(x)
        &=\lim_{\gamma\to 0}\frac{h_i(x)}{\sum_{j\in[n]}h_j(x)}\\
        &=\lim_{\gamma\to 0}\left( 1+ \sum_{j\in[n],j\neq i}\frac{h_j(x)}{h_i(x)}\right)^{-1}\\
        &=\begin{cases}
            0, & x\notin W_i\\
            1, & (x\in W_i) \wedge (x\notin W_j \forall j\in[n]\backslash\{i\})\\
            (\sum_{j\in J}\tfrac{c_j}{c_i})^{-1}, & i\in J\subseteq[n]\colon x\in\bigcap_{j\in J}W_j(x)
        \end{cases}.
    %\end{split}
    \end{align*}
Thus 
\begin{align*}
    H_0p(x)=\lim_{\gamma\to 0}H_\gamma p(x)= \lim_{\gamma\to 0}\gamma S(p*g_\gamma)(x)=-\sumi (x-\mu_i)z_i(x).
\end{align*}
This concludes the proof.
\end{proof}

\begin{proof}[Derivation of equation (\ref{eq:correction_factor})]
    For $S_\theta(\sigma,x)=-\tfrac{x-\mu}{\sigma^2}$,
    we have
    \begin{align*}
        \ddsigma S_\theta(\sigma,x) = \tfrac{2(x-\mu)}{\sigma^3}
    \end{align*}
    and thus (\ref{eq:modified_inference2}) evaluates to
    \begin{align*}
        x_{i+1}
        &= x_{i} - m(t_i)(t_{i+1}-t_i)t_i \big((1+\gamma(t_i))(-\tfrac{x_i-\mu}{t_i^2}) +\tfrac{b\gamma(t_i)}{2t_i}\tfrac{2(x_i-\mu)}{t_i^3})\\
        &= x_{i} - m(t_i)\tfrac{t_i-t_{i+1}}{t_i}(x_i-\mu)\big(1+\gamma(t_i)- \tfrac{b\gamma(t_i)}{t_i^2}\big).
    \end{align*}
    We would like to force this to match the standard score inference step in (\ref{eq:standard_score_Dirac}), for this special case of $S_\theta(\sigma,x)$, which is achieved by choosing 
    \begin{align*}
    m(t_i)=(1+\gamma(t_i)-\tfrac{b\gamma(t_i)}{t_i^2})^{-1}.
    \end{align*}
\end{proof}



\section{Additional Numerical Examples}\label{sub:additionalnumerics}

Additional examples related to the datasets FFHQ and AFHQv2 are shown in Figure~\ref{fig:faces_appendix}. Additional examples related to ImageNet  are shown in Figure~\ref{fig:INET_appendix}.


\begin{figure}
    \centering
    \includegraphics[width=0.99\linewidth]{FFHQ_ab.png}
    \caption{Generated with Algorithm~\ref{alg:main} from the same latent noise sample with $\Delta=0.0001$, $p=8$, and, from left to right, $b\in\{1,2,3,4,5,6,7,8,9,10\}$ as well as, from top to bottom, $a\in\{1,3,5,7\}$. }\label{fig:FFHQ_ab}
\end{figure}


\begin{figure}
    \centering
    \begin{subfigure}{0.42\textwidth}
    \includegraphics[width=\linewidth]{FFHQ_appendix.jpeg}
    \subcaption{FFHQ}
    \end{subfigure}
    \begin{subfigure}{0.42\textwidth}
    \includegraphics[width=\linewidth]{AFHQ_appendix.jpeg} 
    \subcaption{AFHQv2}
    \end{subfigure}
    \caption{In both subfigures, all images in each row starts from the same latent noise sample and the leftmost column uses standard score whereas second to last columns use Algorithm~\ref{alg:main} with $\Delta=0.0001$, $p=8$, $a=2.5$, and, from left to right, $b\in\{2,5,10,20,40,80\}$. The rows are generated from consecutive random seeds.}\label{fig:faces_appendix}
\end{figure}

\begin{figure}
    \centering
    \includegraphics[width=0.86\linewidth]{INET_appendix.jpeg}
    \caption{All images in each row starts from the same latent noise sample and the leftmost column uses standard score whereas second to last columns use Algorithm~\ref{alg:main} with $\Delta=0.001$, $p=12$, $a=4$, and, from left to right, $b\in\{1,2,5,10,15,20,25,30,35,40,50,60,70\}$. The rows are generated from consecutive random seeds.}\label{fig:INET_appendix}
\end{figure}

Figure~\ref{fig:synth_ab_grid} and Figure~\ref{fig:real_ab_grid} show that the samples generated by extended score inference have a strong dependence on the chosen hyperparameters, but a significant emergence of structure can be observed for many different choices. Despite being a simpler problem, the generation seems to be less stable w.r.t.\@ hyperparameter choices for the synthetic data. One reason for this might be that we have a $1$-dimensional manifold of images that is harder to find than the one underlying the EMPIAR-11618 data, which is, of course, not explicitely known but can be assumed to be higher-dimensional. It may also be due to a shorter training duration ($3$mimg compared to $10$mimg), which would be consistent with the generation being much more stable w.r.t.\@ the hyperparameters for FFHQ, AFHQv2, and ImageNet as the networks for those problems have been trained significantly longer. 




\begin{figure}
    \centering
    \includegraphics[width=\linewidth]{synth_ab_grid2.png}
    \caption{Generated with Algorithm~\ref{alg:main} from the same latent noise sample with $\Delta=0.02$, $p=8$, and, from left to right, $b\in\{10,15,20,25,30,35,40,45,50,55,60,65,70,75\}$ as well as, from top to bottom, $a\in\{0.0015,0.002,0.0025,0.003,0.0035,0.004,0.0045,0.005,0.0055\}$.}
    \label{fig:synth_ab_grid}
\end{figure}


\begin{figure}
    \centering
    \includegraphics[width=\linewidth]{real_ab_grid2.jpeg}
    \caption{Generated with Algorithm~\ref{alg:main} from the same latent noise sample with $\Delta=0.001$, $p=8$, and, from left to right, $b\in\{0.25,0.5,1,2,4,8,16,32,64,128,256,512,1024,2048\}$ as well as, from top to bottom, $a\in\{0.0005,0.001,0.002,0.005,0.01,0.02,0.05,0.1,0.2\}$.}
    \label{fig:real_ab_grid}
\end{figure}

In Figure~\ref{fig:synth-pixel}, we show the results with the synthetic dataset with 50\% pixel removal. We see that images generated by standard score inference replicate the corruption present in the training set, while the images generated by extended score inference exhibit all key features of the clean data, while eliminating the corruption due to pixel removal.

\begin{figure}
    \centering
    \includegraphics[width=0.49\linewidth]{rev_synth_PR_modified.png}
    \includegraphics[width=0.49\linewidth]{rev_synth_PR_euler.png}
    
    \caption{Comparison of generated images using a network trained on our synthetic dataset with pixel removal noise, specifically in a given image every pixel is set to $0$ with probability $0.5$. Images generated by standard score inference (right). Images generated by extended score inference (left).}
    \label{fig:synth-pixel}
\end{figure}

In Figure~\ref{fig:paths}, we compare the inference paths for standard and extended score generations: applying truncated sampling to the standard score generation would not be enough to obtain the denoising effect.


\begin{figure}
    \centering
    \includegraphics[trim={0 5.65cm 0 0},clip,width=\linewidth]{rev_path_euler.png}
    
    \vspace{0.2cm}
     
    \includegraphics[trim={0 5.65cm 0 0},clip,width=\linewidth]{rev_path_modified.png}
    \caption{Illustrating the inference path for standard score generation (top) compared to extended score generation (bottom), displaying every second image generated during the inference procedure starting at step $20$.}
    \label{fig:paths}
\end{figure}




Additional images from the EMPIAR-11618 dataset, as well as images generated by standard score inference, are shown in Figure~\ref{fig:real_generated}.

\begin{figure}
    \centering
    \includegraphics[width=0.49\linewidth]{real_data36.jpeg}
    \includegraphics[width=0.49\linewidth]{real_standard.jpeg}
    \caption{Images from the EMPIAR-11618 dataset (left). Images generated by standard score inference (right).}
    \label{fig:real_generated}
\end{figure}


\section{LLM use}

LLMs were used to polish the writing for parts of the text, to suggest related work, and as a coding aid. All those suggestions have only been implemented after thorough manual review. No LLMs were involved in any way in the mathematical derivations. 


\end{document}
